\chapter{Production transverse-$q$ (coupled-$q$) anharmonic transport}
\label{sec:prod_coupledq}

This appendix documents the port of the coupled-$q$ three-phonon
self-energy into the \emph{production} distributed \texttt{quatrex} phonon
solver and the runs that validate and exercise it. The transversely-periodic
physics and the silicon thin-film benchmark are established in the standalone
(dense) solver (\cref{sec:res_qfilm}); the contribution here is software and
verification, not new physics. The full numerical record is saved under
\texttt{phonon/scripts/out/production\_coupled\_q/} (run snapshots,
\texttt{summary.json}) and is reproduced by
\texttt{phonon/scripts/verify/production\_coupled\_q\_summary.py}.

\section{Implementation}

Before this work the production self-energy
(\texttt{src/quatrex/phonon/sse\_phonon\_phonon.py}) raised
\texttt{NotImplementedError} for any \texttt{kpoint\_grid} with a
component $>1$: only the $\Gamma_\perp$-only finite device was supported.
Two changes were made (both with accompanying regression tests; all
$12$ self-energy and $18$ bubble unit tests remain green):

\begin{itemize}
\item \textbf{The coupled-$q$ convolution.} The transversely-periodic
three-phonon self-energy couples the transverse momenta by
crystal-momentum conservation,
$\Sigma(q_\perp,\omega)=\tfrac{1}{N_q}\sum_{q_\perp'}\widetilde\Phi(q_\perp')\,
G(q_\perp')\,G(q_\perp-q_\perp')\,\widetilde\Phi(q_\perp-q_\perp')$, an
extra sum over the internal momentum on top of the energy convolution.
The decisive observation is that in the production data layout only the
energy axis is partitioned across \texttt{comm.stack}; the transverse-$q$
axes of \texttt{global\_stack\_shape} are kept whole and local on every
rank. The momentum sum is therefore \emph{local} --- it required changing
only the inner contraction of the FFT-first bubble, not the
distribution, halo, or Dyson machinery. A guard rejects
\texttt{comm.block.size}${>}1$ together with $k{>}1$ (the spatial-block
halo is not yet momentum-aware), so the film is distributed over the
energy/stack axis.
\item \textbf{Zero-mode projection for $k>1$.} The cell rigid-mode
projector read the dynamical-matrix block and squeezed a single leading
axis, which crashed on the $q$-stacked block of a periodic device; it now
reduces every leading axis to zero, selecting the $q_\perp{=}0$
(\,$\Gamma$-transverse\,) cell matrix as the reference for the rigid/soft
modes. This makes zero-mode projection usable on the film.
\end{itemize}

The $q$-folded vertices are built offline (they do not depend on $G$) and
serialised through a small loader
(\texttt{src/quatrex/phonon/qfold.py}); the production solver only reads
arrays and carries no dependency on the standalone science code.

\section{Verification}

The port was checked to be \emph{bit-exact} against independent
references before any physics run (\cref{tab:prod_validation}). These are
the gate: the production self-energy reproduces an independently written
einsum oracle of the coupled-$q$ convolution to machine precision, the
production-assembled film Hamiltonian $H(q_\perp)$ reproduces the dense
\texttt{get\_btd\_blocks} to $10^{-10}$, and the real-space
dynamical-matrix file used by the production input path round-trips
against the dense $H(q_\perp)$ exactly.

\begin{table}[h]
\centering
\caption{Bit-exact verification of the coupled-$q$ port.}
\label{tab:prod_validation}
\begin{tabular}{lc}
\hline
Check & relative error \\
\hline
coupled-$q$ self-energy vs.\ independent einsum oracle & $6.4\times10^{-16}$ \\
production $H(q_\perp)$ vs.\ dense \texttt{get\_btd\_blocks} & $1.1\times10^{-10}$ \\
dynamical-matrix transverse-IDFT round-trip & $5.0\times10^{-16}$ \\
\hline
\end{tabular}
\end{table}

\paragraph{Input pipeline.} The production \texttt{construct\_from\_unit\_cell}
path folds transverse Bloch phases over a Monkhorst--Pack mesh. Because
that mesh is half-shifted, $q=(k+\tfrac12)/n_k-\tfrac12+\text{shift}$,
while the standalone solver uses the $\Gamma$-centred $q=k/n_k$, the two
were reconciled by setting \texttt{kpoint\_shift}$=\tfrac12-\tfrac{1}{2n_k}$
per transverse axis. The real-space \texttt{dynamical\_matrix.mat} is then
the exact transverse inverse-DFT of the dense $H(q_\perp)$ over that mesh,
which is why the round-trip in \cref{tab:prod_validation} is exact by
construction. The film inputs use the bulk-silicon force constants of
\cref{sec:res_qfilm} with transport along a primitive axis.

\section{Silicon thin film: transverse-mesh sweep}
\label{sec:prod_film}

The production film was run with the recipe that proved numerically
stable (see \cref{sec:prod_open}): broadening $\eta=0.4$\,THz, $\eta_{\rm
obc}=0.1$, retarded method \texttt{half}, and zero-mode projection of the
three $q_\perp{=}0$ acoustic translations (which sit at $0$\,THz and carry
no heat). The transverse mesh was swept at fixed thickness ($3$ primitive
layers) to test momentum convergence, with one thicker point
(\cref{tab:prod_film}, \cref{fig:prod_qfilm_qconv}).

\begin{table}[h]
\centering
\caption{Production silicon film, coupled-$q$ self-energy, $\eta=0.4$\,THz
with zero-mode projection. $G$ are the $q$-summed lead heat currents
(arbitrary, common units); $G/N_q$ is the per-momentum value. The lead
non-conservation is $|J_0-J_{-1}|/|J_0|$. The dense benchmark
(\cref{sec:res_qfilm}) at a comparable setting is shown for scale.}
\label{tab:prod_film}
\begin{tabular}{llccc}
\hline
$n_k$ & $n_{\rm slabs}$ & $G_{\rm anh}/G_{\rm ball}$ & $G_{\rm ball}/N_q$ & lead non-cons. \\
\hline
$3$ & $3$ & $0.861$ & $3.29$ & $9.7\%$ \\
$5$ & $3$ & $0.871$ & $3.02$ & $9.5\%$ \\
$7$ & $3$ & $0.867$ & $3.05$ & $9.6\%$ \\
$5$ & $5$ & $0.846$ & --- & $15.1\%$ \\
\hline
\multicolumn{2}{l}{(no zero-mode, $n_k{=}3$, $n_{\rm slabs}{=}3$)} & $0.733$ & --- & $22\%$ \\
\hline
\end{tabular}
\end{table}

\begin{figure}[h]
\centering
\includegraphics[width=0.6\linewidth]{transport_sweeps/prod_qfilm_qconv.pdf}
\caption{Production silicon film coupled-$q$ transport vs.\ transverse mesh
($n_{\rm slabs}{=}3$). The per-momentum ballistic and anharmonic
conductances settle by $n_k{=}5$ and the anharmonic ratio is
momentum-converged at $G_{\rm anh}/G_{\rm ball}\approx0.87$.}
\label{fig:prod_qfilm_qconv}
\end{figure}

Two facts: (i) the per-momentum conductance and the anharmonic ratio are
\emph{momentum-converged} --- $G_{\rm ball}/N_q$ settles by $n_k{=}5$ and
the ratio is $0.861/0.871/0.867$ across $n_k=3/5/7$ --- so the coupled-$q$
machinery behaves correctly across the full mesh; (ii) zero-mode
projection of the $0$\,THz translations roughly halves the
non-conservation ($22\%\to10\%$) and raises the ratio
($0.733\to0.861$ at $n_k{=}3$), the same soft-mode effect seen on the
wires. The residual $\sim10\%$ non-conservation does \emph{not} shrink
with $n_k$ (\cref{fig:prod_qfilm_conservation}); it is therefore not a
momentum-undersampling effect and is treated as an open issue in
\cref{sec:prod_open}.

\begin{figure}[h]
\centering
\includegraphics[width=0.6\linewidth]{transport_sweeps/prod_qfilm_conservation.pdf}
\caption{Lead-to-lead heat non-conservation of the production film. It is
flat in the transverse mesh and is reduced, but not removed, by zero-mode
projection. The dense solver reaches $\sim10^{-5}$ on the same physical
system at small broadening; the gap is the open issue of
\cref{sec:prod_open}.}
\label{fig:prod_qfilm_conservation}
\end{figure}

\section{Carbon nanotube: length ladder}
\label{sec:prod_cnt}

To exercise the $k{=}1$ path at production length scales (where the dense
multi-cell solver is prohibitively expensive, \cref{sec:res_cnt}), the
$(3,3)$ carbon nanotube was run as a length ladder $L=2,4,6$ transport
cells at $\eta=0.45$\,THz on a $121$-point grid with $50$ iterations
(\cref{fig:prod_cnt_ladder}). The anharmonic self-consistency on this
soft-twist tube is metastable: the iteration does not settle to a fixed
point but passes through a conserving iterate that the heat-flow criterion
captures (the final iterate diverges and is discarded). The captured
best-iterate ratios are $G_{\rm anh}/G_{\rm ball}=0.896$ ($L{=}2$,
conservation $1.0\%$) and $0.911$ ($L{=}4$, conservation $3.5\%$); the
$L{=}6$ point is not shown. The large-broadening value is
weaker (closer to one) than the dense small-$\eta$ ladder of
\cref{sec:res_cnt}, consistent with the under-scattering discussed in
\cref{sec:prod_open}.

\begin{figure}[h]
\centering
\includegraphics[width=0.55\linewidth]{transport_sweeps/prod_cnt_ladder.pdf}
\caption{Production CNT(3,3) anharmonic ratio (best conserving iterate) vs.\
device length at $\eta=0.45$\,THz; per-point lead-to-lead conservation
annotated.}
\label{fig:prod_cnt_ladder}
\end{figure}

\section{Cross-check against the dense solver}
\label{sec:prod_dense}

To separate port correctness from broadening/grid effects, the dense
\texttt{si\_film} solver was run at the \emph{same} transverse mesh
($n_k{=}3$, $n_{\rm slabs}{=}3$) at two broadenings
(\cref{tab:prod_dense}). At the production-accessible broadening
($\eta_w\approx0.4$) the dense solver also fails to conserve
($17\%$) and gives a ratio above one (the documented single-cell
broadening artefact, \cref{sec:res_approx}); only at small broadening
($\eta_w\approx0.05$) does the dense ratio settle to $0.766$. The
production ($0.733$--$0.861$) and the dense both live in the unreliable
large-broadening regime here, which is why their ratios differ; neither
large-broadening ratio is the physical answer. Notably the dense solver
also conserves only to $13\%$ at $n_k{=}3$, far from the $\sim10^{-5}$ it
reaches at the converged $n_k{\geq}8$ mesh of \cref{sec:res_qfilm} ---
i.e.\ the coarse $n_k{=}3$ mesh used for the matched cross-check is itself
inadequate for conservation in \emph{both} solvers.

\begin{table}[h]
\centering
\caption{Dense \texttt{si\_film} cross-check, $n_k{=}3$, $n_{\rm slabs}{=}3$,
transport along the same axis as the production run.}
\label{tab:prod_dense}
\begin{tabular}{lccc}
\hline
broadening & $G_{\rm anh}/G_{\rm ball}$ & conservation \\
\hline
$\eta_w\approx0.4$ (matched to production) & $1.022$ & $17\%$ \\
$\eta_w\approx0.05$ (small) & $0.766$ & $13\%$ \\
\hline
\end{tabular}
\end{table}

\section{Open issues}
\label{sec:prod_open}

\paragraph{Heat-flow conservation ($\sim10\%$) --- open, under
investigation.} The production film conserves the lead-to-lead energy
current only to $\sim10\%$ (with zero-mode projection; $22\%$ without),
independent of the transverse mesh, against the dense solver's
$\sim10^{-5}$ on the same physical system. A code-level candidate has been
identified but \emph{not yet confirmed} as the dominant cause. The
FFT-first bubble is \emph{index-based}: it forms
$\Sigma(\omega_m)=\sum_n G_n\,G_{m-n}$ (the frequency grid enters only
through the spacing $d\omega$ and the retarded Hilbert transform; the
bubble call uses \texttt{zero\_freq\_idx=None}). This is exact only if the
grid begins at $\omega_0=0$, so that difference frequencies
$\omega_m-\omega_n=(m-n)\,d\omega$ land on grid points. The film runs used
$\omega_0=0.05$\,THz $=0.38\,d\omega$ (non-commensurate), whereas the CNT
runs used $\omega_0=0.01$\,THz $=0.02\,d\omega$; the observed
non-conservation tracks $\omega_0/d\omega$ (film $\sim10\%$, CNT
$\sim1\%$). Whether aligning the grid to $\omega_0=0$ (with proper
DC-point handling, as the dense solver does) and/or a larger/finer grid
recovers the dense conservation is the subject of the follow-up
investigation; it is not asserted here as established.

\paragraph{Small-broadening instability --- needs a solver-level fix.} The
production film is numerically stable only for $\eta\gtrsim0.4$\,THz; at
$\eta\leq0.2$ the Dyson solve $[(\omega+i\eta)^2 I - D - \Sigma^R]$ becomes
singular and the block recursion returns NaN, where the dense solver's
dense factorisation stays finite. The config-level remedies were
exhausted without success: retarded method \texttt{fft},
large $\eta_{\rm obc}$ ($0.4$--$0.8$ with bulk $\eta\in[0.15,0.2]$), and
zero-mode projection all still produced NaN at $\eta\leq0.2$. A
solver-level regularisation of the block recursion is therefore required
to reach the small-broadening regime the dense solver uses. This is a
pre-existing fragility of the production phonon SCBA (the same family as
the soft-mode behaviour of \cref{sec:res_nw}), not specific to the
coupled-$q$ port.

\paragraph{Large-broadening under-scattering.} As a consequence of the
above, the production-accessible ($\eta\approx0.4$) anharmonic reduction
is systematically weaker than the dense small-broadening result --- film
$0.87$ vs.\ dense $0.77$, CNT $0.90$ vs.\ dense $0.70$--$0.81$
(\cref{sec:res_cnt}). The same weakening appears in both the $k{>}1$ film
and the $k{=}1$ tube, consistent with it being a broadening effect rather
than anything specific to the momentum coupling.
